\documentclass[11pt]{article}
\usepackage[margin=1in]{geometry}
\usepackage{amsmath}
\usepackage{amssymb}
\usepackage{booktabs}
\usepackage{listings}
\usepackage{xcolor}
\usepackage{hyperref}

\lstset{basicstyle=\ttfamily\small,breaklines=true,frame=single,
        backgroundcolor=\color{gray!8},columns=fullflexible}
\newcommand{\code}[1]{\texttt{#1}}

\title{Is the reversed-wind adjoint problem caused by integrating\\
       backward in time? Transpose versus inverse}
\author{}
\date{}

\begin{document}
\maketitle

\noindent\emph{This note answers a specific conceptual question that arises when
reading the two reversed-wind advection-adjoint variants
(\code{ADV\_ADJ\_METHOD = reverse\_wind\_pure} vs.\ \code{reverse\_wind}): are the
cube-seam ``corrections'' needed because the adjoint integrates in reverse time,
or would negating the input wind and running an ordinary forward simulation be
enough? The short answer is that backward-in-time integration is necessary and
correct and is \emph{not} the source of the corrections; the corrections exist
because the adjoint requires the \emph{transpose} of the transport operator,
which is a different object from the \emph{inverse} that a reversed-wind forward
run produces. They coincide in the smooth interior and differ at the asymmetric
grid seams, and the corrections are exactly that difference.}

%======================================================================
\section{Two operations both march ``backward''}

The adjoint method inherently marches backward in time: computing
$\partial J/\partial(\text{state})$ by the chain rule assembles the sensitivity
from the last time step to the first, so the saved trajectory is always visited
in reverse order. There is no forward-in-time gradient adjoint. Hence reverse-time
integration is not optional, and it is not something we could remove.

But ``march backward'' can mean two mathematically distinct things, and this is
the crux:
\begin{itemize}
  \item \textbf{Time-reversal / inverse.} Run the dynamics backward to
  \emph{undo} the transport. This is what ``negate the wind and integrate
  forward'' produces: an approximate $A^{-1}$.
  \item \textbf{Transpose (the adjoint).} Apply $A^\top$ of each step, in reverse
  order. This propagates sensitivities and satisfies the dot-product identity
  $\langle A x, y\rangle = \langle x, A^\top y\rangle$.
\end{itemize}
Both visit the steps in reverse order. The transpose additionally \emph{flips the
operator}, not merely the calendar. For a single forward step represented by a
matrix $A$, integrating backward as an inverse applies $A^{-1}$, whereas the
adjoint applies $A^\top$, and in general
\begin{equation}
  A^{-1} \neq A^\top .
\end{equation}
The reversed-wind trick substitutes an inverse-like forward pass for the
transpose. That substitution is exact only where $A$ is self-adjoint
(symmetric); elsewhere it is merely an approximation, and the corrections repair
the rows where it fails.

%======================================================================
\section{Where inverse and transpose agree, and where they split}

\subsection{Interior: reversed-wind forward \emph{equals} the transpose}

Consider one-dimensional linear advection with constant velocity and first-order
upwind. Forward ($u>0$, Courant number $c=u\,\Delta t/\Delta x$, $0<c<1$):
\begin{equation}
  q_i^{\text{new}} = (1-c)\,q_i + c\,q_{i-1},
\end{equation}
i.e.\ the matrix row $A_{i,i}=1-c$, $A_{i,i-1}=c$. Its transpose is
\begin{equation}
  (A^\top q)_i = (1-c)\,q_i + c\,q_{i+1}.
\end{equation}
Negating the wind ($u<0$, upwind now from the right) and running the same forward
scheme gives
\begin{equation}
  q_i^{\text{new}} = (1-c)\,q_i + c\,q_{i+1},
\end{equation}
which is \emph{identical} to $(A^\top q)_i$. So in the constant-coefficient
interior, reversed-wind-forward already \emph{is} the transpose, for free and
with no corrections. This is the sense in which ``just reverse the wind'' is
correct, and why \code{reverse\_wind\_pure} is not nonsense in cell interiors.

\subsection{Seam/boundary: they differ by the cross-seam term}

Put this on a bounded domain, cells $1..N$, and inspect the last row.
\begin{itemize}
  \item \textbf{Transpose row $N$:} $(A^\top q)_N = (1-c)\,q_N + c\,q_{N+1}$, but
  cell $N+1$ lies \emph{outside} the domain; the forward operator never had an
  interior entry $A_{N+1,N}$, so that term does not exist. The transpose row is
  simply $(1-c)\,q_N$.
  \item \textbf{Reversed-wind forward at $N$:} a valid forward run reconstructs
  using a ghost cell, computing $(1-c)\,q_N + c\,q_{N+1}^{\text{ghost}}$, where on
  the cubed sphere $q_{N+1}^{\text{ghost}}$ is filled from the neighbouring face
  by \code{copy\_corners}.
\end{itemize}
The two differ \emph{by exactly the ghost term}. The reversed-wind forward run
includes a cross-seam contribution the transpose does not have---and it must
include it, because pulling that ghost is what makes it a well-posed forward
integration. This is why negating the input wind cannot avoid the problem: no
matter which way the wind points, a correct forward run at a boundary
reconstructs from the off-face ghost. \emph{The property that makes it a good
forward simulation is precisely the property that makes it a wrong transpose at
the seam.}

The two \code{reverse\_wind} corrections are exactly this repair:
\begin{itemize}
  \item the \textbf{interior-only edge reconstruction} deletes the spurious ghost
  term ($c\,q_{N+1}^{\text{ghost}}$) so the row matches the transpose;
  \item \code{copy\_corners\_ad} re-routes the corner sensitivity to where the
  transpose sends it (accumulated back onto the source real cells) instead of
  scattering it outward to a ghost.
\end{itemize}

%======================================================================
\section{The causal chain}

The role of reverse-time is upstream of the corrections, not identical to them:
\begin{quote}
adjoint method $\Rightarrow$ \textbf{must march backward in time} $\Rightarrow$
sensitivity propagates \emph{opposite} to mass $\Rightarrow$ so we flip the wind
to mimic the transpose $\Rightarrow$ flipping the wind reproduces the transpose
\emph{only where the operator is symmetric} $\Rightarrow$ the seams are
asymmetric $\Rightarrow$ corrections are needed.
\end{quote}
Reverse-time is the \emph{reason we flip the wind}, but the \emph{bug} is one step
further down: flipping the wind is only an approximate transpose, and it
over-reaches at the asymmetric boundary rows.

%======================================================================
\section{The decisive test: \code{exact\_linear} also runs backward in time}

The cleanest way to isolate the variable is to compare against the exact adjoint.
\code{exact\_linear} \emph{also} integrates in reverse time: it steps through the
saved trajectory backward and even reverses the sub-step order---the forward does
horizontal advection then vertical remap, while the exact adjoint does the
vertical-remap transpose first, then the horizontal transpose. Same reverse-time
direction as \code{reverse\_wind}.

Yet \code{exact\_linear} requires \emph{none} of the \code{copy\_corners\_ad} or
interior-only edge-stencil corrections, because it transposes each operator
\emph{exactly}: the corner handling and boundary stencils are transposed
faithfully inside its own adjoint routines rather than being approximated by a
reversed-wind forward pass. Same time direction, exact spatial transpose, no seam
corrections.

Holding reverse-time fixed, the corrections appear only for the \emph{approximate}
(reversed-wind) transpose and never for the \emph{exact} one. Therefore the
corrections are about how faithfully the \emph{spatial} operator is transposed,
not about integrating backward in time.

\begin{center}
\begin{tabular}{lccc}
\toprule
variant & reverse time? & spatial transpose & seam corrections? \\
\midrule
\code{reverse\_wind\_pure} & yes & reversed-wind (approx.) & none (hence sign error) \\
\code{reverse\_wind}       & yes & reversed-wind $+$ seam fixes & yes \\
\code{exact\_linear}       & yes & exact transpose & not needed \\
\bottomrule
\end{tabular}
\end{center}

%======================================================================
\section{One-line answer}

Integrating in reverse time is \emph{why} the wind gets flipped, but it is not the
problem. The problem is that a reversed-wind forward run is the operator's
\emph{inverse}-like partner, whereas the adjoint needs its \emph{transpose}. The
two agree in the self-adjoint interior and disagree at the asymmetric grid seams,
and the corrections are exactly that disagreement. Run the exact transpose
(\code{exact\_linear}) and the same reverse-time integration needs no corrections
at all.

\end{document}
